\begin{table}[p]\centering
\caption{Directional targets: step size and certified oracle error}\label{tab:r12-fd}
\setlength{\tabcolsep}{3.5pt}\begin{tabular}{rrrr}
\toprule
Step size & \shortstack{Additional endpoint\\error allowance} & \shortstack{Maximum measured\\direction error} & \shortstack{Maximum certified\\target-error bound} \\
\midrule
$1/80$ & 0e+00 & 5.58e-07 & 0.00208 \\
$1/80$ & 1e-07 & 5.58e-07 & 0.00209 \\
$1/80$ & 1e-05 & 5.58e-07 & 0.00288 \\
$1/40$ & 0e+00 & 8.89e-06 & 0.00417 \\
$1/40$ & 1e-07 & 8.89e-06 & 0.00417 \\
$1/40$ & 1e-05 & 8.89e-06 & 0.00457 \\
$1/20$ & 0e+00 & 3.45e-05 & 0.00833 \\
$1/20$ & 1e-07 & 3.45e-05 & 0.00834 \\
$1/20$ & 1e-05 & 3.45e-05 & 0.00853 \\
\bottomrule
\end{tabular}
\begin{minipage}{.98\textwidth}\small
Twelve independently generated training-sized problems. Endpoint values are re-solved and rationally certified at every step size. The extra error column is an explicit sensitivity allowance added to each endpoint certificate, not a claim that fresh noisy labels were simulated. Measured error compares with the separately solved central actor and is numerical; the displayed bound includes oracle gaps and the finite-difference bias bound.
\end{minipage}
\end{table}
